\documentclass[wabun2025.tex]{subfiles}
\begin{document}
\section{Discussion}
\label{sec:analysis}
\subsection{Configuration of the Simulator Circuit}
\label{sec:circuit-structure-analysis}
\subsubsection{Comparison of Resource Consumption Between the Two Circuit Methods}
As shown in Figure \ref{fig:hamiltonian-gatecount}, the required number of
qubits for both the arithmetic circuit and the UCRZ circuit increases linearly
with the number of coordinate qubits. However, the arithmetic operation has a
larger constant term and slope. On the other hand, looking at the number of
2-qubit gates as a measure of circuit depth, the number of UCRZ gates increases
exponentially with respect to the input (dimensionality $\times$ number of
coordinate qubits), whereas the arithmetic circuit shows only a moderate
increase. However, in small-scale models, the smallness of the constant term for
the UCRZ gate is dominant. Therefore, a trade-off was confirmed: direct state
preparation using UCRZ gates is advantageous for small-scale, low-dimensional
systems, while arithmetic operations will be superior in terms of depth when
performing large-scale, high-resolution simulations in the future.

\subsubsection{Utility of UCRZ Gates in Quantum Emulators}
In quantum emulators on classical computers, the simulatable state vector space
is exponentially limited by the hardware's memory (RAM) capacity (typically
around 20 to 30 qubits). Between the two circuit methods examined in this study,
the approach of using UCRZ gates to suppress the consumption of ancilla qubits
is extremely useful for algorithm verification, as it circumvents this severe
qubit constraint and enables Hamiltonian simulations of higher dimensions (such
as the 2D hydrogen atom model in this study) to be executed on the emulator.

\subsubsection{Prospects for Implementation Costs Anticipating Application to FTQC Devices}
While any operational gate can be executed at a uniform cost on an emulator, the
situation is different for actual fault-tolerant quantum computers (FTQC). In
particular, promising error correction schemes like the Surface code require
enormous physical resources, such as magic state distillation, to execute
non-Clifford gates compared to Clifford gates. The CC-NOT gates in arithmetic
circuits and arbitrary Rz rotation gates in UCRZ circuits fall precisely into
this category. Ultimately, the bottleneck lies in how many T gates these can be
approximated and decomposed into (T-count / T-depth). In particular, the Rz
rotation in the UCRZ gate significantly fluctuates in T-gate consumption
depending on the required accuracy. Therefore, for application to real devices,
a quantitative comparison of the gate implementation costs of both methods,
taking into account the overhead of error correction, will be a future
challenge.

\subsection{Handling the Pole of the Coulomb Potential}
\label{sec:coulomb-potential-handling-analysis}
\subsubsection{Reason Why the Effective Radius is $\deltar$/4}
As shown in Figure \ref{fig:Ha1-vs-delta1}, the optimal value for the effective
radius $\reff=\deltar/a$ used in the local average potential was obtained at
$a=4$ for the 2D model. That is, setting $\reff=\deltar/4$ most effectively
regularized the divergence caused by the pole (singularity) of the Coulomb
potential. This result was intriguing and differed from our initial intuition.

Initially, we focused on the energy expectation value $\braket{\psi|H|\psi}$
within a circular region of radius $\deltar/2$ centered on the singularity, and
determined $\reff$ such that the integral value using the analytical solution
matched the approximation formula using the representative value. The $\reff$
derived by this method includes higher-order terms of the grid spacing $\deltar$
and approaches $\reff=\deltar/4$ in the limit of $\deltar\rightarrow0$. Assuming
that these higher-order terms could not be ignored under conditions of coarse
coordinate resolution (e.g., $\nb=6$), we performed time evolution simulations
using $\reff$ including the higher-order terms. However, comparative
verification revealed that the system's energy value approached the analytical
solution more closely when using $\deltar/4$, where the higher-order terms are
truncated.

To investigate the cause of this, we re-evaluated $\reff$ so that the geometric
regional average of the simple potential function $V(r)=-1/r$, without weighting
by the wave function, matched the representative value. This directly derived
$\reff=\deltar/4$ without higher-order terms. This fact suggests that in a
finitely discretized grid space, adopting a simple mean field of the microscopic
region represented by the lattice point serving as the singularity as the
potential yields more accurate energy, rather than assuming higher-order changes
of the wave function within the grid.

\subsubsection{Validity of the Approximation Within a Circular Region}
To obtain a local average centered on a lattice point, one method is to
integrate over a square region with side $\deltar$ represented by that lattice
point. However, in this study, to obtain an analytically closed-form expression,
the calculation was performed in a circular region of radius $\deltar/2$
centered on the singularity (lattice point). The integral value of the Coulomb
potential is dominated by the contribution near the origin, which is the
singularity, and is isotropic. Therefore, the circular region is considered to
capture the essential spatial scaling better. The fact that the energy value in
the numerical experiment was optimal at the value ($a=4$) obtained from this
approximation supports the validity of this approximation.

\subsubsection{Reason Why Truncating Higher-Order Terms Improves Accuracy}
As mentioned above, higher-order terms of $\deltar$ appear in a rigorous
evaluation using $\braket{\psi|H|\psi}$, but applying this to the simulation
ironically decreased the accuracy. This is presumed to be because, on a coarse
grid with a finite grid spacing $\deltar$, the higher-order spatial changes
(curvature) of the wave function around the lattice point cannot be fully
expressed. As a result, it became clear that adopting a simple geometric average
(truncating higher-order terms), which does not depend on minute shape changes
of the wave function exceeding the grid resolution, provides the most stable
regularization parameter.

\subsubsection{Extension to 3D Models and General Versatility}
\label{subsec:discussion-3d-extension}

To verify whether the local average potential method as a technique to avoid the
singularity of the Coulomb potential is an essential property independent of
dimensions (the general versatility of this method), we evaluated the extension
of this method to a 3D model.

In 3D space, a minute cubic grid (side $\deltar$) containing the singularity is
approximated by a sphere of radius $\deltar/2$ centered on the singularity, and
the spatial average of the Coulomb potential $-1/r$ inside it is calculated.
Assuming the volume of the sphere of radius $R = \deltar/2$ is $V =
\frac{4}{3}\pi R^3$, the average potential inside the sphere is obtained
analytically as follows.
\begin{equation}
    \frac{1}{V} \int_{0}^{R} -\frac{1}{r} 4\pi r^2 dr 
    = \frac{1}{\frac{4}{3}\pi R^3} \cdot 4\pi
      \left[ -\frac{1}{2}r^2 \right]_{0}^{R} 
    = -\frac{3}{2R} = -\frac{3}{\deltar}=-\frac{1}{\reff}
\end{equation}
Therefore, the effective radius for avoiding divergence at the singularity in
the 3D model is derived as $\reff = \deltar/3$.

To verify this theoretical prediction, numerical calculations for a 3D hydrogen
atom model were performed. Since it is difficult to secure a practical number of
qubits for a 3D model using currently available quantum circuit emulators, a
classical simulation program that performs time evolution operations equivalent
to quantum circuits was used for this evaluation.

The results are shown in Figure \ref{fig:Ha1-vs-delta1-3d} of Appendix
\ref{sec:appendix_a_3d-evaluation}. Setting the effective radius as
$r_{\mathrm{eff}} = \deltar/a$, the coefficient $a$ was varied from $1$ to $5$,
and the energy expectation value of the ground state was evaluated after
executing time evolution for a sufficient period under conditions of spatial
resolution $\nb \in \{6, 7, 8\}$. As a result, similar to the 2D model, the
energy expectation value matched the analytical ground state energy of the 3D
hydrogen atom model ($-0.5\ \mathrm{a.u.}$) when $a=3$ (i.e., $r_{\mathrm{eff}}
= \deltar/3$) under all $\nb$ conditions.

This fact indicates that the approach of the local average potential, which
determines the effective radius using the geometric spatial average of a minute
region containing the singularity, is effective regardless of the
dimensionality. This suggests that when fault-tolerant quantum computers (FTQC)
with a sufficient number of qubits are realized in the future, and
first-quantization simulations of practical 3D molecular models are performed,
the Coulomb singularity can be regularized with high accuracy using the simple
and clear guideline of $\deltar/3$.

\subsection{Simulation Parameters for Minimizing Error}
\label{sec:parameter-selection-for-minimum-error-analysis}
\subsubsection{Dynamic Range of Simulation at a Given $\nb$}
In general, the target system handled in a simulation is represented by a linear
combination of multiple eigenstates with different quantum numbers, and these
need to be represented simultaneously under a single spatial grid spacing
$\deltar$.

Based on the trade-off analysis between the spatial discretization error and the
domain truncation error (cutoff error) with respect to $\deltar$ shown in Figure
\ref{fig:dq-tradeoff}, the ground state, which has the narrowest wave function
spread, requires the smallest $\deltar$. When attempting to represent up to the
highest possible excited state simultaneously with the ground state, it is
desirable to set the value of $\deltar$ as large as possible within a range that
does not hinder the representation of the ground state. At that time, the
maximum quantum number that can be represented without being affected by the
cutoff error is determined, and the range from the ground state to that maximum
quantum number becomes the representable range of this simulation, that is, the
"dynamic range." This serves as a powerful guideline for estimating the
representable state space in advance when executing simulations with limited
quantum resources.

\subsubsection{Increase in $\nb$ and the Plateau of Optimal Parameters}
Looking at Figure \ref{fig:dq-tradeoff}, it was confirmed that as $\nb$
increases, the plateau (flat region) of $\deltar$ where both the spatial
discretization error and the cutoff error are negligible expands. This means
that as the number of available qubits increases, the system becomes more robust
against parameter setting errors. In an environment where the hardware is
immature and $\nb$ is small, strict tuning of $\deltar$ is required, but if
$\nb$ increases in the future, it is expected that the simulation accuracy will
become more tolerant to parameter selection.

\subsubsection{Comparison Between Worst-Case Trotter Error Estimates and Calculated Errors}
Conventional Trotter error evaluations tend to require an extremely small
$\deltat$ because they assume the worst case based on the operator norm
\cite{Burgarth2024.PhysRevResearch.6.043155}. However, in realistic systems
where energy is concentrated in eigenstates near the ground state,
state-dependent error evaluation is more appropriate. The results of this
experiment show that results close to the accuracy limit for a given $\nb$ can
be obtained by using a $\deltat$ that satisfies the limit of phase change based
on the sampling theorem (half-period condition), which is coarser than the
formal requirements of the Trotter error. This makes it possible to avoid
increasing the depth of the quantum circuit caused by unnecessarily fine time
steps.

\subsubsection{Trade-off Between Spatial Error and Time Error}
As shown in Figure \ref{fig:trotter-error-t1-part1}, a phenomenon was observed
where the reduction of the system's energy error saturated even when $\deltat$
was made finer than a certain threshold. This can be interpreted as the dominant
error factor of the overall simulation shifting from the time error (Trotter
error or sampling error) to the spatial error (mainly spatial discretization
error) caused by the coarseness of the spatial grid $\deltar$. That is, there is
a limit to the spatial accuracy that can be achieved with a certain combination
of $\nb$ and $\deltar$, and making $\deltat$ disproportionately fine (i.e.,
making the quantum circuit deeper) leads to a waste of quantum resources. This
analysis provides an integrated procedure for determining the spatial resolution
$\deltar$ and the time step $\deltat$ in a well-balanced manner for a given
number of qubits $\nb$.

\end{document}